\subsection{代数中范数与迹的性质}

命题3. 设 A 是一个具有有限基的 K-代数。元素 \(a \in \mathbf{A}\) 可逆的必要且充分条件是它的 \(\mathbf{N}_{\mathbf{A} / \mathbf{K}}(a)\) 在 \(\mathbf{K}\) 中可逆。

若 \(a\) 具有逆元 \(a'\) （位于 \(\mathbf{A}\) 中），则

\[
\mathrm{N} _ {\mathbf {A} / \mathbf {K}} (a) \mathrm{N} _ {\mathbf {A} / \mathbf {K}} (a ^ {\prime}) = \mathrm{N} _ {\mathbf {A} / \mathbf {K}} (a a ^ {\prime}) = \mathrm{N} _ {\mathbf {A} / \mathbf {K}} (1) = 1
\]

由第 1 号公式 (3)。反之，若 \(\mathbf{N}_{\mathbf{A} / \mathbb{K}}(a)\) 可逆，则自同态 \(h: x \mapsto ax\) 是双射（§ 8，第 2 号，定理 1）。于是存在 \(a' \in A\) 使得 \(aa' = 1\) ；于是 \(h(a'a - 1) = aa'a - a = (aa' - 1)a = 0\) ，从而 \(aa' = 1\) ，因为 \(h\) 是单射。因此 \(a'\) 是 \(a\) 的逆元。

命题4. 设 A 是一个具有有限基的 K-代数。对所有 \(a \in \mathbf{A}\) ， \(\operatorname{Pc}_{\mathbf{A} / \mathbf{K}}(a; a) = 0\) 。

这直接由凯莱-哈密顿定理（§ 8，第 11 号，命题 20）得出。

命题5. 设 A 是一个 K-代数，m 是 A 的一个双边理想。假设 \(A_{0} = A/m\) 是维数为 n 的自由 K-模，存在整数 r \textgreater{} 0 使得 \(m^{r} = \{0\}\) ，并且 \(m^{i-1}/m^{i}\) 是自由 \(A_{0}\) -模，其维数为 \(s_{i}\) （ \(1 \leqslant i \leqslant r\) ）。设 \(s = s_{1} + \cdots + s_{r}\) ，且对所有 \(a \in A\) ，设 \(a_{0}\) 表示 a 模 m 所在的类。则 A 是维数为 ns 的自由 K-模，且对所有 \(a \in A\) ，

\[
\operatorname{Tr} _ {\mathbf {A}} (a) = s. \operatorname{Tr} _ {\mathbf {A} _ {0}} (a _ {0}), \quad \mathrm{N} _ {\mathbf {A}} (a) = (\mathrm{N} _ {\mathbf {A} _ {0}} (a _ {0})) ^ {s}\tag{23}
\]

\[
\mathrm{Pc} _ {\mathbf {A}} (a; \mathbf {X}) = (\mathrm{Pc} _ {\mathbf {A} _ {0}} (a _ {0}; \mathbf {X})) ^ {s}.
\]

根据 II，§1，第 13 号，命题 25， \(\mathfrak{m}^{i - 1} / \mathfrak{m}^{i}\) 是维数为 \(ns_i\) 的自由 K-模。因此，第 2 号中的命题 1 可以应用于 \(\mathbf{P}_i = \mathfrak{m}^{i - 1} / \mathfrak{m}^{i}\) ；

这首先表明 A 是一个维数为 \(n(s_{1} + \cdots + s_{r}) = ns\) 的自由 K-模。此外，该假设蕴含 A-模 \(P_{i}\) 同构于 \(s_{i}\) 个子模的直和，这些子模均同构于 A-模 \(A_{0}\) ；由第 2 号命题 1，因此 \(\mathrm{N}_{\mathrm{P}_{i}}(a) = \mathrm{N}_{\mathrm{A}_{0}}(a)^{s_{i}}\) ；最终因此

\[
\mathrm{N} _ {\mathbf {A}} (a) = \mathrm{N} _ {\mathbf {A} _ {0}} (a) ^ {s}.
\]

在此公式中， \(\mathrm{N}_{\mathbf{A}_{0}}(a)\) 是通过将 \(A_{0}\) 视为左 A-模来定义的，它等于 K-线性映射 \(x \mapsto ax\) 在 \(A_{0}\) 到自身上的行列式；但由于 \(ax = a_{0}x\) 对 \(x \in A_{0}\) 成立，故 \(\mathrm{N}_{\mathbf{A}_{0}}(a) = \mathrm{N}_{\mathbf{A}_{0}}(a_{0})\) ，这就完成了范数公式 (23) 的证明。另外两个公式可类似地证明。

命题6. 设 A 是 K 上的一个交换代数，且在 K 上具有有限基，V 是一个 A-模，且在 A 上具有有限基。则 V 在 K 上具有有限基，并且对 V 的每个 A-自同态 u，若 \(u_{\mathbf{K}}\) 是将u视为V的K自同态的映射，

\[
\operatorname{Tr} (u _ {\mathbb {K}}) = \operatorname{Tr} _ {\mathbf {A} / \mathbb {K}} (\operatorname{Tr} (u)), \quad \det (u _ {\mathbb {K}}) = \mathrm{N} _ {\mathbf {A} / \mathbb {K}} (\det (u))\tag{24}
\]

\[
\operatorname{Pc} (u _ {\mathbf {K}}; \mathbf {X}) = \mathrm{N} _ {\mathbf {A} [ \mathbf {X} ] / \mathbf {K} [ \mathbf {X} ]} (\operatorname{Pc} (u; \mathbf {X})).
\]

设 \((a_{i})_{1\leqslant i\leqslant m}\) 是 A 在 K 上的一个基， \((e_{j})_{1\leqslant j\leqslant n}\) 是 V 在 A 上的一个基；则 \((a_{i}e_{j})\) 是 V 在 K 上的一个基（II，§1，第13号，命题25）。另一方面，将第二个公式应用于自同态 \(\mathbf{X} - \bar{\boldsymbol{u}}\) （即 A{[}X{]}-模 A{[}X{]} \(\otimes_{\mathbf{A}}\mathrm{V}\) 的自同态）（§ 8，第 10 号），即可导出公式 (24) 中的第三个。因此，只需证明 (24) 中的前两个公式即可。我们首先建立以下引理：

引理1. 设 \(\mathbf{X}_{ij}\) ( \(1 \leqslant i \leqslant n\) , \(1 \leqslant j \leqslant n\) ) 是 \(n^2\) 个未定元， \(X\) 是元素为 \((\mathbf{X}_{ij})\) 的 \(n\) 阶方阵， \(\mathbf{D}(\mathbf{X}_{11}, \ldots, \mathbf{X}_{nn}) \in \mathbf{Z}[\mathbf{X}_{11}, \ldots, \mathbf{X}_{nn}]\) 是行列式 \(\det(X)\) 。另一方面，设 A 为交换环， \(\mathbf{M}_{ij}\) ( \(1 \leqslant i \leqslant n\) , \(1 \leqslant j \leqslant n\) ) 是 A 上的 \(n^2\) 个 \(m\) 阶矩阵，它们两两可交换，M 是 A 上的 \(mn\) 阶方阵，它可以表示为矩阵的分块方阵（II，§10，第7号）

\[
M = \left( \begin{array}{c c c c} M _ {1 1} & M _ {1 2} & \ldots & M _ {1 n} \\ M _ {2 1} & M _ {2 2} & \ldots & M _ {2 n} \\ \cdot & \cdot & \cdot & \cdot \\ M _ {n 1} & M _ {n 2} & \ldots & M _ {n n} \end{array} \right)
\]

那么M的行列式等于方阵的行列式

\[
\mathrm{D} (M _ {1 1}, \dots , M _ {n n})
\]

该方阵的阶为 m。

我们对 \(n\) 进行归纳，情形 \(n = 0\) 和 \(n = 1\) 是平凡的。

设 \(Z\) 为一个新的未定元， \(N_{ij}\) 表示矩阵 \(M_{ij} + \delta_{ij}ZI_m\) ( \(\delta_{ij}\) 为克罗内克指标)。若 \(D^{ij}(X_{11},\ldots,X_{nn})\) 是 \(X_{ij}\) 在矩阵 \(X\) （§ 8，第 6 号）中的代数余子式，则

\[
\mathrm{X} _ {j l} \mathrm{D} ^ {k l} (\mathrm{X} _ {1 1}, \dots , \mathrm{X} _ {n n}) = \delta_ {j k} \mathrm{D} (\mathrm{X} _ {1 1}, \dots , \mathrm{X} _ {n n})\tag{25}
\]

（§ 8，第 6 号，公式 (28)）。我们记 \(N_{ij}' = \mathbf{D}^{ij}(N_{11}, \ldots, N_{nn})\) ，这是一个 \(m\) 阶方阵，位于 \(\mathrm{A}[Z]\) 上；并考虑乘积 \(N.U\) ，其中

\[
U = \left( \begin{array}{c c c c} N _ {1 1} ^ {\prime} & 0 & \ldots & 0 \\ N _ {1 2} ^ {\prime} & I _ {m} & \ldots & 0 \\ \cdot & \cdot & \cdot & \cdot \\ N _ {1 n} ^ {\prime} & 0 & \ldots & I _ {m} \end{array} \right),
\]

\[
N = \left( \begin{array}{c c c c} N _ {1 1} & N _ {1 2} & \ldots & N _ {1 n} \\ N _ {2 1} & N _ {2 2} & \ldots & N _ {2 n} \\ \cdot & \cdot & \cdot & \cdot \\ N _ {n 1} & N _ {n 2} & \ldots & N _ {n n} \end{array} \right)
\]

按分块进行此乘积（II，§ 10，第 5 号）并利用公式 (25)，我们得到

\[
N. U = \left( \begin{array}{c c c c} P & N _ {1 2} & \ldots & N _ {1 n} \\ 0 & N _ {2 2} & \ldots & N _ {2 n} \\ \cdot & \cdot & \cdot & \cdot \\ 0 & N _ {n 2} & \ldots & N _ {n n} \end{array} \right)
\]

其中我们记作 \(P = \mathrm{D}(N_{11},\ldots ,N_{nn})\) 。设

\[
Q = \left( \begin{array}{c c c c} N _ {2 2} & \ldots & N _ {2 n} \\ \cdot & \cdot & \cdot & \cdot \\ N _ {n 2} & \ldots & N _ {n n} \end{array} \right)
\]

，这是一个 \(m(n - 1)\) 阶矩阵；则（§8，第6号，公式 (31)）(det \(N\) ) ( \(\det U\) ) = ( \(\det P\) ) ( \(\det Q\) ) 且 \(\det U = \det N_{11}'\) 。但根据归纳假设， \(\det Q = \det(\mathrm{D}^{11}(N_{11},\ldots ,N_{nn})) = \det N_{11}'\) ，并且根据 \(N_{ij}\) 的定义，显然 \(\det Q\) 是 A{[}Z{]} 中一个次数为 \(m(n - 1)\) 的多项式，其 \(Z^{m(n - 1)}\) 项的系数为 1；由此立即得出 \(\det Q\) 在分次代数 A{[}Z{]} 中不是零因子。因此我们得出结论： \(\det N = \det(\mathrm{D}(N_{11},\ldots ,N_{nn}))\) 在 A{[}Z{]} 中成立；如果我们在这些多项式中将 Z 替换为 0，那么 \(\det M = \det(\mathrm{D}(M_{11},\ldots ,M_{nn}))\) 。

证明了这一引理之后，K-模 V 是诸 K-模 \(Ae_{j}\) ( \(1 \leqslant j \leqslant n\) ) 的直和；我们记 \(u(e_{j}) = \sum_{k=1}^{n} c_{jk} e_{k}\) 。对每个元素 \(xe_{j} \in Ae_{j}\) ，其中 \(x \in A\) ， \(u(xe_{j})\) 在 \(Ae_{k}\) 中的分量是 \(xc_{jk} e_{k}\) ；由此可知， \(u_{k}\) 关于 K-模 V 的基 \((a_{i} e_{j})\) 的矩阵可以表示为分块方阵 \((M_{jk})\) 的形式，其中 \(M_{jk}\) 是 K-线性映射 \(xe_{j} \mapsto xc_{jk} e_{k}\) （它把 \(Ae_{j}\) 映到 \(Ae_{k}\) ）关于这两个 K-模的基 \((a_{i} e_{j})_{1 \leqslant i \leqslant m}\) 和 \((a_{i} e_{k})_{1 \leqslant i \leqslant m}\) （II，§ 10，第 5 号）的矩阵。若对所有 \(t \in A\) ， \(M(t)\) 表示 A 在 K 上的基 \((a_{i})_{1 \leqslant i \leqslant m}\) 下 A 的自同态 \(x \mapsto xt\) 的矩阵，则 \(M_{jk} = M(c_{jk})\) ；由于 \(t \mapsto M(t)\) 是一个环同态，矩阵 \(M_{jk}\) 两两可交换。则

\[
\det u _ {\mathrm{K}} = \det (\mathrm{D} (M _ {1 1}, \dots , M _ {n n}))
\]

由引理1。但由于 \(t \mapsto M(t)\) 是一个环同态， \(\mathrm{D}(M_{11},\ldots ,M_{nn})\) 是 A 的 K-自同态 \(x \mapsto x.\det (c_{jk})\) 关于基 \((a_t)\) 的矩阵；根据定义，其行列式因此为 \(\mathrm{N}_{\mathbf{A} / \mathbf{K}}(\det (u))\) ，这证明了 (24) 中的第二个公式。另一方面，

\[
\operatorname{Tr} \left(u _ {\mathrm{K}}\right) = \sum_ {j = 1} ^ {n} \operatorname{Tr} \left(M _ {j j}\right) = \sum_ {j = 1} ^ {n} \operatorname{Tr} _ {\mathrm{A} / \mathrm{K}} \left(c _ {j j}\right) = \operatorname{Tr} _ {\mathrm{A} / \mathrm{K}} \left(\sum_ {j = 1} ^ {n} c _ {j j}\right) = \operatorname{Tr} _ {\mathrm{A} / \mathrm{K}} (\operatorname{Tr} (u))
\]

并且命题6的证明完成。

推论。设A是交换K-代数，在K上具有有限基，B是A-代数，在A上具有有限基。则B在K上具有有限基，并且，对所有 \(b \in B\) （“传递性公式”）

\[
\operatorname{Tr} _ {\mathbf {B} / \mathbf {K}} (b) = \operatorname{Tr} _ {\mathbf {A} / \mathbf {K}} (\operatorname{Tr} _ {\mathbf {B} / \mathbf {A}} (b)), \quad \mathrm{N} _ {\mathbf {B} / \mathbf {K}} (b) = \mathrm{N} _ {\mathbf {A} / \mathbf {K}} (\mathrm{N} _ {\mathbf {B} / \mathbf {A}} (b))\tag{26}
\]

\[
\mathrm{Pc} _ {\mathbf {B} / \mathbf {K}} (b; \mathbf {X}) = \mathrm{N} _ {\mathbf {A} [ \mathbf {X} ] / \mathbf {K} [ \mathbf {X} ]} (\mathrm{Pc} _ {\mathbf {B} / \mathbf {A}} (b; \mathbf {X})).
\]

这直接从命题 6 得出，取 \(\mathbf{V} = \mathbf{B}\) 和 \(u(x) = bx\) 。

注记。假设 K 到 A 的同态 \(\lambda \mapsto \lambda, 1\) 是单射，并将 K 与其在 A 中的像等同；假设 A 作为 K-模具有有限基 \((e_i)_{1 \leqslant i \leqslant n}\) 。设 \(s\) 是 A 的一个自同构，使得 \(s(\mathbf{K}) = \mathbf{K}\) 。设 \(a\) 是 A 的一个元素；则通过转移结构

\begin{enumerate}
\def\labelenumi{(\arabic{enumi})}
\setcounter{enumi}{26}
\tightlist
\item
\end{enumerate}

\[
\operatorname{Tr} _ {\mathbf {A} / \mathbf {K}} (s (a)) = s (\operatorname{Tr} _ {\mathbf {A} / \mathbf {K}} (a))\tag{28}
\]

\[
\mathrm{N} _ {\mathbf {A} / \mathbf {K}} (s (a)) = s (\mathrm{N} _ {\mathbf {A} / \mathbf {K}} (a)).
\]

同时考虑 A 的一个导子 D（§ 10，第 2 号），使得 D(K) ⊂ K，并记 \(D (e _ {i}) = \sum_ {j = 1} ^ {n} e _ {j} \mu_ {j i}\) ，其中 \(\mu_ {j i} \in \mathbf {K}\) ；记

\[
a e _ {i} = \sum_ {j = 1} ^ {n} e _ {j} \lambda_ {j i} \quad \text { with } \quad \lambda_ {j i} \in \mathbf {K}.
\]

那么

\[
\mathrm{D} (a) e _ {i} + a \mathrm{D} (e _ {i}) = \mathrm{D} (a e _ {i}) = \sum_ {j = 1} ^ {n} (\mathrm{D} (e _ {j}) \lambda_ {j i} + e _ {j} \mathrm{D} (\lambda_ {j i})).
\]

由此可得

\[
\mathbf {D} (a) e _ {i} = \sum_ {j = 1} ^ {n} e _ {j} \nu_ {j i}
\]

其中 \(\nu_{ji} = \mathrm{D}(\lambda_{ji}) + \sum_{k=1}^{n} (\mu_{jk}\lambda_{ki} - \lambda_{jk}\mu_{ki})\) 。由于 \(\sum_{i,k} (\mu_{ik}\lambda_{ki} - \lambda_{ik}\mu_{ki}) = 0\) ，因此

 \(\mathrm{Tr}_{\mathbf{A} / \mathbb{K}}(\mathbf{D}(a)) = \sum_{i = 1}^{n}\mathbf{D}(\lambda_{ii})\) ，换言之

\[
\operatorname{Tr} _ {\mathbf {A} / \mathbf {K}} (\mathbf {D} (a)) = \mathbf {D} (\operatorname{Tr} _ {\mathbf {A} / \mathbf {K}} (a)). *\tag{29}
\]

